\subsection{Environmental impacts}

The environmental impacts of \acp{bss} have been a central part of their justification~\cite{Shaheen2010, Fishman2016, Midgley2011}. These impacts include transport emissions, energy use, material consumption, and end-of-life waste and recycling. While early academic discussions often focused on \ac{ghg} reductions from substituting car trips, recent research has broadened the perspective to include the full life cycle of shared bicycles, docking infrastructure, batteries, and logistic operations.

    \subsubsection{Transport emissions reductions derived from modal substitution} 

Transport emissions encompass two related but distinct pollutant categories: \acp{ghg} such as carbon dioxide (\ce{CO2}), methane (\ce{CH4}), and nitrous oxide (\ce{N2O}), which accumulate in the atmosphere, and local air pollutants such as nitrogen oxides (\ce{NOx}), hydrocarbon (\ce{HC}), particulate matter (\ce{PM_{2.5}} and \ce{PM_{10}}), and carbon monoxide (\ce{CO}), which are short-lived and primarily affect urban air quality. Both categories are produced by motorized transport modes, but their implications differ: \ac{ghg} reductions contribute to climate mitigation, while pollutant reductions directly improve street-level air quality.

Early environmental assessments often reported large climate benefits, typically by estimating the total distance travelled by bike-sharing and then attributing these kilometres to alternative modes using assumed substitution patterns. Then, emission savings were calculated by applying mode-specific emission factors to the replaced kilometres and subtracting the negligible emissions of cycling.

In Barcelona, for example, \textit{Bicing} was estimated to avoid more than 9,000 tonnes of \ce{CO2} in 2011, but assuming that 90\% of users were new cyclists and that all of their trips would otherwise have been made by car~\cite{Rojas-Rueda2011}. Later work in Shanghai estimated savings of 25,000 tonnes of \ce{CO2} in 2016. However, the study assumed that all trips shorter than 1~km replaced walking, while longer trips replaced taxi, and the annual citywide results were extrapolated from a single month of data by adjusting for the sample ratio (56.62\%), the company’s market share (56.56\%), and scaling up to twelve months~\cite{Zhang2018}. At the same time, Zhang~\cite{Zhang2018} acknowledged that the environmental benefits of bike-sharing are inherently localized in time and space, since they are directly correlated with patterns of bicycle use rather than being uniform across the city and year. Other studies reported even larger figures. Qiu and He~\cite{Qiu2018}, for instance, estimated a reduction of 616,036 tonnes of \ce{CO2} in 2020 compared to the 2015 value, assuming that 75\% of bike-sharing kilometres were replacing car trips. An analysis of New York’s \textit{Citi Bike} marked some progress in this field by mapping trips onto the street network rather than using straight-line distances, but its estimates of 30,070 tonnes of \ce{CO2} reductions between 2014 and 2017 still rested on deterministic distance-based thresholds borrowed from other research~\cite{Chen2022c}.

Taken together, these approaches tend to overstate benefits mainly because they make strong assumptions about modal substitution (more details on modal substitution in Section~\ref{sec:mobility_and_transport_impacts}). To address these shortcomings, studies refined how modal substitution is modelled. Kou et al.~\cite{Kou2020} developed a trip-level probabilistic model that combined observed bike-sharing distances with city-specific travel survey data and life-cycle emission factors, assigning each trip to a most likely replaced mode. Using this method, they estimated annual reductions across eight U.S. cities in 2016 ranging from 41 to 5,417 tonnes of \ce{CO2}, or roughly 0.3--0.6~kg per trip. Similarly, Saltykova~\cite{Saltykova2022}, analyzing a dockless \ac{bss} in Chengdu, used traveller survey responses to quantify modal substitution. They found that 39\% of bike-sharing trips replaced bus journeys, 13.5\% replaced subway, and only 10.9\% replaced car trips. With these figures, the avoided emissions fell markedly compared to a “cars-only” counterfactual—demonstrating how omitting \ac{pt} substitution inflates climate benefits by a factor of two or more.

A more refined approach was proposed by Li et al.~\cite{Liu2021}, who developed a high resolution discrete-choice framework matching over 23 million dockless bike-sharing trips in Shanghai to realistic modal alternatives using travel time, distance, and cost information from online routing services. Their results showed that most trips substituted for walking (30.7\%) or \ac{pt} (50\%), with car replacement below 2\%, leading to estimated annual savings of 116,000 tonnes of \ce{CO2}. This progression from simple distance rules to survey-informed probabilities, and finally to discrete-choice models, illustrates a gradual shift toward more behaviourally grounded and context-sensitive estimates of the environmental benefits.

Research has also estimated reductions in local air pollutants following the same logic as with \ac{ghg} emissions, although such assessments are far less common. For example, in Shanghai, bike-sharing was estimated to cut \ce{NOx} emissions by 64 tonnes~\cite{Zhang2018}; while in Beijing, \ce{SO2}, \ce{CO}, and \ce{NO2} emissions were reduced by 22.5, 58.64, and 1,586.7 tonnes, respectively, and \ce{PM_{2.5}} and \ce{PM_{10}} by 10.3 tonnes~\cite{Qiu2018}.


    \subsubsection[LCA for BSSs]{\ac{lca} for \acp{bss}}

The environmental impacts resulting from the operation of \acp{bss} are not limited to avoided transport emissions due to increased bike usage. For example, Fishman, Washington, and Haworth~\cite{Fishman2014} showed that London’s bike-sharing scheme in 2012 actually produced higher overall motor vehicle usage because of the rebalancing operations needed to keep the service running. Therefore, to fully capture the impact resulting from bike usage, additional factors must be taken into account. The manufacturing of bicycles, docks, and batteries, as well as the maintenance, logistics operations, and end-of-life handling when hardware is no longer usable, all require material- and energy-intensive processes, contributing to embodied emissions.

\paragraph{LCA methodology overview\\}

To capture these impacts, studies typically employ \acp{lca}, usually based on the ISO 14,040 and 14,044 standards~\cite{LCA2006a, LCA2006b}. These standards define how to account for the burdens associated with manufacturing, operation, maintenance, and end-of-life management (Figure~\ref{fig:LCA_bss}). Importantly, full \acp{lca} should also incorporate the benefits from modal substitution, so that avoided emissions from displaced car or \ac{pt} trips are evaluated alongside the embodied burdens of system operation.

    \begin{figure}[htbp!]
        \centering
        \includegraphics[width=1\textwidth, trim={0cm 0.3cm 0cm 1cm}, clip]{0_plots/figure_LCA.pdf}
        \caption{Aspects that should be considered in a \acs{bss} \ac{lca}. Figure based on Luo et al.~\cite{Luo2019}, adapted and elaborated by the author.}
        \label{fig:LCA_bss}
    \end{figure}

Within the \ac{lca} framework, the life-cycle inventory phase quantifies all relevant material and energy flows entering and leaving the system. These inventories constitute the empirical foundation for the subsequent impact assessment stage, where the quantified flows are translated into environmental indicators. To perform this translation, researchers apply established methods such as ReCiPe 2008~\cite{Bonilla-Alicea2020}, ReCiPe 2016~\cite{Felipe-Falgas2022}, or GREET~\cite{Cazzola2020, Krauss2022}. These approaches allow impacts to be expressed either at the midpoint level, which reports specific categories such as climate change or particulate matter formation with relatively low uncertainty, or at the endpoint level, which aggregates them into broader areas of protection such as human health, ecosystems, and resources, with higher uncertainty but greater communicability. Earlier studies of shared micromobility often limited their scope to \ac{ghg} emissions, while more recent \acp{lca} studies present both midpoint and endpoint indicators~\cite{Sun2022, Chen2024, Calan2024}.

By applying these methodologies, researchers can compute comparable indicators across systems. At the midpoint level, results are typically reported per service delivered, in grams of \ce{CO2}-eq per passenger-kilometer (pkm) of bike travel~\cite{Luo2019, Chen2024, Bonilla-Alicea2020}. The term \ce{CO2}-eq (“carbon dioxide equivalent”) expresses the combined impact of different \acp{ghg}, such as \ce{CO2}, \ce{CH4}, and \ce{N2O}, by converting them into the amount of \ce{CO2} that would cause the same warming effect, based on their global warming potentials. This metric enables comparison across systems, but also with other transport modes~\cite{Sun2022, Krauss2022}. Reported values for complete \acp{bss} \acp{lca} range from roughly 30 to 150~g~\ce{CO2}-eq/pkm~\cite{Zhu2023, Felipe-Falgas2022, Cazzola2020, Sun2022, Calan2024}. This wide variation reflects differences in system type, usage and operational practices, and whether all life-cycle phases are included in the analysis.

\paragraph{Manufacturing phase impacts\\}

The first step of an \ac{lca} is assessing the manufacturing phase, which covers the extraction of raw materials and production of bicycles and docking stations, and batteries where applicable. This is modelled by compiling life-cycle inventories of material inputs—such as aluminium and steel for frames, lithium and cobalt for batteries, plastics for saddles and tyres, and electronics for smart bikes or stations—and linking them with emission factors from databases like Ecoinvent~\cite{Felipe-Falgas2022} or CPCD~\cite{Zhu2023}. 

Manufacturing often emerges as the largest single contributor, accounting for around 50\% of total life-cycle emissions in docked systems~\cite{Calan2024, Sun2022}. When comparing systems with \acp{e-b}, the production of the hardware accounts 61\% of the life-cycle emissions for docked systems and 52\% for dockless systems~\cite{Luo2019}. Moreover, Bonilla-Alicea~\cite{Bonilla-Alicea2020} quantified manufacturing burdens and found that dockless systems' bicycles with on-board electronics had approximately 3.7 times the climate change impact of docked bicycles, and about 2.7 times their total endpoint impacts. The increase was largely driven by the electronic subsystems (solar panels, printed circuit boards, communication modules, and batteries), highlighting that these components are candidates for technological improvement in terms of environmental impact. On the other hand, \acp{e-b} add further burdens from battery and motor production, with batteries alone contributing 20–30\% of embodied \acp{ghg}~\cite{Zhu2023}. Material choices also play a decisive role. Earlier \acp{lca} of private bicycles showed that aluminium frames carried higher production impacts~\cite{Dave2010, Leunberger2010, Luna2016} than alternatives such as steel or bamboo~\cite{Agyekum2017}, a finding echoed in shared fleets where heavier, vandal-resistant frames and integrated electronics further raise embodied emissions.

\paragraph{Usage phase impacts\\}

The usage phase of \acp{bss} is not limited to the act of cycling itself, which is practically emission-free, but also involves several supporting operations. \acp{lca} typically account for rebalancing, battery charging, maintenance, and modal substitution savings. 

Among these, \textbf{rebalancing} consistently emerges as the dominant contributor, responsible for 36\% of total life-cycle emissions in docked systems and up to 73\% in dockless ones~\cite{Luo2019}, which corresponded to average values around 118~g~\ce{CO2}-eq/pkm for dockless systems and 65--70~g~\ce{CO2}-eq/pkm for docked ones~\cite{Luo2019}.

The high burden of rebalancing arises because it requires motorized service vehicles to redistribute bicycles across the city. Its impact depends on several factors: the number of bicycles in the system, overall system size, the types of vehicles used for redistribution, and how often rebalancing occurs~\cite{Bonilla-Alicea2020}. To quantify this burden, \acp{lca} model the distance travelled by rebalancing vehicles relative to the service delivered, often expressed as service-vehicle kilometers per bicycle-kilometer. Estimates range from as low as 0.012 km of van travel per bike-km in docked systems~\cite{Felipe-Falgas2022} to 0.10–0.26 km in dockless fleets~\cite{Sun2022, Cazzola2020}. Some systems, for instance, rely on conventional vans with intensities around 300~g~\ce{CO2}-eq/km~\cite{Chen2020LifeCycle, Hollingsworth2019}, while others have begun experimenting with electric vans or cargo bikes for battery collection and charging, significantly reducing the associated emissions~\cite{Felipe-Falgas2022, Calan2024}. More broadly, improving rebalancing efficiency has become a research field of its own, with numerous studies developing optimization models and operational strategies to minimize environmental and financial costs.

Moreover, \textbf{battery charging} should also be considered. As Sun et al.~\cite{Sun2022} emphasize, electrically assisted systems incur burdens not only from the collection and redistribution of bicycles, but also from the energy required to recharge batteries. The direct electricity demand of shared \acp{e-b} is relatively modest—around 0.014--0.017~kWh per km, or roughly 1.5~kWh per 100~km~\cite{Sun2022}—yet its impact depends strongly on the carbon intensity of the grid, ranging from about 10--20~g~\ce{CO2}-eq/km in low-carbon contexts to more than 70~g~\ce{CO2}-eq/km in coal-dominated regions~\cite{Zhu2023, Calan2024}. Although small compared to the emissions from rebalancing logistics, it still should be considered as a part of the use-phase footprint.

Another component of the use phase is \textbf{maintenance operations}, which include routine activities such as component replacement, pumping, lubricating, and cleaning~\cite{Zhu2023}. Because operator records are rarely available~\cite{Calan2024}, \acp{lca} typically rely on interviews~\cite{Zhu2023}, technical manuals, or web sources~\cite{Bonilla-Alicea2020} to estimate service frequencies and part lifetimes. As a result, maintenance is often underrepresented in system assessments. When explicitly included, maintenance is usually found to contribute around 10–15\% of total life-cycle \acp{ghg} in conventional docked systems~\cite{Sun2022, Chen2024, deBortoli2021Environmental}. These shares can rise in dockless fleets, where higher rates of wear and vandalism accelerate the need for component replacement~\cite{Luo2019}.

Overall, the relative importance of maintenance depends strongly on assumptions about vehicle lifetimes and component durability. Shorter lifespans for bicycles and batteries increase the share of maintenance, whereas strategies that extend durability and improve repairability can substantially mitigate impacts~\cite{Zhu2023, Calan2024}.      

Beyond operational burdens, the use phase in most \acp{lca} also credits \textbf{avoided emissions from modal substitution}. In practice, this means assigning each bike-sharing trip a counterfactual mode (e.g., car, \ac{pt}, walking) and subtracting the corresponding life-cycle emissions that would have been produced otherwise. Recent \acp{lca} operationalize this with survey-based mode shares or trip data: in Washington, D.C., Chen et al.~\cite{Chen2024} incorporate mode-shift information to compute a net balance where avoided emissions are weighed against manufacturing and operations. In Barcelona, Felipe-Falgas et al.~\cite{Felipe-Falgas2022} use an \ac{lca} inventory with self-reported substitution and show that shared \acp{e-b} can increase \acp{ghg} when they primarily replace \ac{pt} or walking rather than cars. At a broader scale, Sun et al.~\cite{Sun2022} formalize the accounting as the difference between the bike-sharing emissions and a weighted mix of displaced modes (plus new trips), reporting results in g~\ce{CO2}-eq/pkm.

Several syntheses reinforce this integrated “burdens minus benefits” framing. Krauss et al.~\cite{Krauss2022} adapted mode-shift survey data and \acp{lca} factors to compute the net sustainability impact of shared micromobility across six cities, showing that outcomes depend critically on the share of car and ride-hailing trips displaced versus low-carbon modes. Methodological guidance from the ITF/OECD likewise treats modal shift as essential to comparability, recommending city-specific substitution profiles and transparent functional units (pkm) when reporting net outcomes~\cite{Cazzola2020}. Together, these studies demonstrate that whether bike sharing reduces \acp{ghg} in the use phase, depends less on cycling’s direct emissions (which are near zero) and more on what trips it replaces and how operations are managed.


\paragraph{End-of-life phase impacts\\}

The final step of an \ac{lca} involves evaluating the end-of-life phase, which covers the collection and subsequent recycling or disposal of retired bicycles, batteries, docks, and electronic components. Similar to the maintenance phase, this stage is often underrepresented in \acp{lca} due to data limitations~\cite{Calan2024}. Notably, Krauss et al.~\cite{Krauss2022} is among the few studies to include end-of-life impacts; their findings indicate that these processes represent only a small share of total life-cycle emissions.

To address the lack of empirical data, some researchers have adopted informed assumptions regarding disposal pathways. For instance, Krauss et al.~\cite{Krauss2022} and Sun et al.~\cite{Sun2023} assume that 80–90\% of metal components are recycled, while the remainder are either incinerated or sent to landfill. Comparative studies have found that incineration typically yields slightly lower net \ac{ghg} emissions compared to landfill, though in both cases, the emissions from disposal are generally only about one-fifth those associated with manufacturing~\cite{Zhu2023, Chen2024}. As a result, the overall influence of the disposal route on total \ac{ghg} emissions is relatively minor.

However, real-world recycling outcomes can vary widely, depending on factors such as local regulations and the availability of recycling technology~\cite{Sun2022, Calan2024}. Recovery rates for different materials range from 30\% to 90\%. While metals like aluminium and steel have high recovery potential, plastics, electronics, and especially batteries pose greater recycling challenges and are often handled as waste. For example, in Washington, D.C., Chen~\cite{Chen2024} showed that shifting from landfill to waste processing reduced the end-of-life carbon footprint of a shared bike by about 6~kg~\ce{CO2}-eq, and systematic metal recycling can provide even greater emission offsets.

The treatment of batteries is especially important in this context. Batteries already account for significant impacts during production, and managing their end-of-life poses additional environmental risks due to hazardous materials and recycling inefficiencies~\cite{Felipe-Falgas2022, Calan2024}. Without robust battery recycling schemes in place, these disposal challenges could negate some of the emission savings achieved through modal substitution.

Beyond the direct recycling and disposal process, other end-of-life factors, such as fleet attrition and oversupply, can further exacerbate environmental impacts. For example, when around 20\% of bikes are lost each year to vandalism or theft, overall fleet emissions may rise by as much as 25\%, since new units are required to replace those prematurely retired~\cite{Calan2024}. A dramatic demonstration of this challenge was seen in Chinese cities during the 2017–2018 dockless bike boom, when “bike-sharing graveyards” (Fig.~\ref{fig:bss_graveyard}) emerged as excess and abandoned bikes accumulated in massive lots~\cite{Sun2022}. In these situations, low utilization prevented manufacturing emissions from being distributed across enough trips, and inadequate recycling meant much of the material value was lost~\cite{Luo2019, Sun2022}.

In summary, while the end-of-life phase typically contributes only a small percentage of total \acp{ghg}, its significance comes from the opportunity it provides to close resource loops, recover valuable materials, and prevent hazardous waste accumulation. Therefore, from a sustainability standpoint, implementing robust recycling systems, promoting modular product design, and ensuring appropriately sized fleets are crucial strategies for improving environmental outcomes at the end of a \acp{bss}’s life cycle~\cite{Krauss2022, Calan2024}.

\begin{figure}    
    \centering
    \includegraphics[width=0.8\textwidth]{0_plots/graveyard.jpg}
    \caption{Bike-sharing graveyard in Xiamen, China. Sourced from \textit{Amusing Planet}~\cite{AmusingPlanet2018}}
    \label{fig:bss_graveyard}
\end{figure}

\paragraph{Synthesis\\}

Overall, the environmental performance of \acp{bss} depends strongly on system design, operational practices, and local mobility patterns. While shared bicycles can generate meaningful reductions in transport emissions, these benefits materialize only when car substitution is substantial and when operational burdens—particularly rebalancing and battery logistics—are minimized. At the same time, manufacturing and end-of-life stages introduce non-negligible embodied impacts that must be accounted for when evaluating system sustainability. The evidence therefore suggests that \acp{bss} are not inherently low-carbon by default, but can achieve positive environmental outcomes when implemented as part of a well-designed, efficiently operated, and right-sized multimodal transport system.
